% !TeX root = ../navier-stokes-manuscript.tex
% ============================================================
% 2.4 The Spatial Column Inside Critical Height
% ============================================================

\subsection{The Spatial Column Inside Critical Height}\label{sub:TheSpatialColumnInsideCriticalHeight}

The mixed jet is too large to reveal at once how temporal escalation reaches into space.
Critical height supplies a scalar through which one complete spatial column can be followed.
For \(s\ge0\), set
\[
  E_s(t)
  :=
  \frac12\norm{\Lambda^su(t)}_2^2,
\]
so that
\[
  H(t)=E_{1/2}(t).
\]
Keeping \(s\) free before returning to \(1/2\) exposes the mechanism that repeats at every height.

Apply \(\Lambda^s\) to the pressure-complete momentum equation and pair it in \(L^2\) with \(\Lambda^su\).
The temporal term gives
\[
  \pair{\Lambda^su}{\Lambda^s\partial_tu}
  =
  \frac{d}{dt}\frac12\norm{\Lambda^su}_2^2
  =
  E_s'.
\]
Rapid decay and commutation with spatial derivatives give
\[
  \nu\pair{\Lambda^su}{\Lambda^s\Delta u}
  =
  -\nu\norm{\Lambda^{s+1}u}_2^2
  =
  -2\nu E_{s+1}.
\]
The pressure pairing vanishes only after incompressibility and integration by parts have done their work:
\[
  \pair{\Lambda^su}{\Lambda^s\nabla p}
  =
  -\pair{\nabla\cdot\Lambda^su}{\Lambda^sp}
  =
  0.
\]
This global cancellation does not remove pressure from the local equation or from the mixed jet that generated the balance.
Keep the signed nonlinear transfer in its pressure-complete form,
\[
  \mathcal P_s(t)
  :=
  -\pair{\Lambda^su}
  {\Lambda^s\bigl((u\cdot\nabla)u+\nabla p\bigr)}.
\]
The exact balance is
\[
  E_s'
  =
  \mathcal P_s-2\nu E_{s+1}.
\]
One time response at spatial height \(s\) has opened the velocity energy one full Sobolev step higher.
It has also kept the signed transfer created by the lower mixed rows.

At critical height, the first row reads
\[
  H'
  =
  \mathcal P_{1/2}-2\nu E_{3/2}.
\]
This is already the first spatial ascent, from the scale-invariant \(E_{1/2}\) to \(E_{3/2}\).
It is not a bound for \(E_{3/2}\), because the transfer and dissipative terms can cancel inside the complete value of \(H'\).
The whole row has to be differentiated before its next spatial consequence can be seen.

The balance one step higher is
\[
  E_{3/2}'
  =
  \mathcal P_{3/2}-2\nu E_{5/2}.
\]
Substituting it into the derivative of the first critical row gives
\begin{align*}
  H''
  &=
  \mathcal P_{1/2}'-2\nu E_{3/2}'
  \\
  &=
  \mathcal P_{1/2}'
  -2\nu\mathcal P_{3/2}
  +(2\nu)^2E_{5/2}.
\end{align*}
The highest visible spatial energy has climbed again, and the transfer at \(3/2\) has joined the transfer already present at \(1/2\).
Differentiating once more uses
\[
  E_{5/2}'
  =
  \mathcal P_{5/2}-2\nu E_{7/2}
\]
and yields
\[
  H'''
  =
  \mathcal P_{1/2}''
  -2\nu\mathcal P_{3/2}'
  +(2\nu)^2\mathcal P_{5/2}
  -(2\nu)^3E_{7/2}.
\]
One more response gives
\[
  H^{(4)}
  =
  \mathcal P_{1/2}'''
  -2\nu\mathcal P_{3/2}''
  +(2\nu)^2\mathcal P_{5/2}'
  -(2\nu)^3\mathcal P_{7/2}
  +(2\nu)^4E_{9/2}.
\]
Each temporal derivative reaches one spatial step higher, but it carries every transfer encountered on the way.

The pattern is triangular rather than diagonal.
For every integer \(n\ge1\),
\[
  H^{(n)}
  =
  (-2\nu)^nE_{n+1/2}
  +
  \sum_{j=0}^{n-1}
  (-2\nu)^j
  \partial_t^{\,n-1-j}\mathcal P_{j+1/2}.
\]
To pass from row \(n\) to row \(n+1\), differentiate the existing transfer terms and replace
\[
  E_{n+1/2}'
  =
  \mathcal P_{n+1/2}-2\nu E_{n+3/2}.
\]
The replacement creates the next transfer and the next highest spatial energy with the required coefficient.
This proves the identity by induction, but the identity still supplies no estimate by itself.
Even if \(H^{(n)}\) is known, \(E_{n+1/2}\) cannot be isolated without controlling the entire signed transfer sum.
The highest energy appearing in a row is a coordinate discovered inside that row, not a quantity bounded by its appearance.

The homogeneous scale made the critical dilation exact, while continuation is usually read in the inhomogeneous scale.
For \(s>0\),
\[
  \norm{u(t)}_{H^s}^2
  \simeq
  \norm{u(t)}_2^2+\norm{\Lambda^su(t)}_2^2
  =
  \norm{u(t)}_2^2+2E_s(t).
\]
The unforced energy identity supplies the missing low-frequency part:
\[
  \frac12\norm{u(t)}_2^2
  +
  \nu\int_0^t\norm{\nabla u(\tau)}_2^2\,d\tau
  =
  \frac12\norm{u_0}_2^2.
\]
Thus control of \(E_s(t)\) supplies the high-frequency part of an \(H^s\) bound, and an \(L^1\)-in-time bound for \(E_s\) supplies the high-frequency part of an \(L^2_tH^s_x\) bound on each finite interval.

Critical height first exposed the column because \(1/2\) is fixed by scaling and the first derivative immediately reveals \(3/2\).
The proof implication at the endpoint runs in the opposite logical direction.
Once the equation-generated velocity jet is placed in every finite Sobolev row on the whole terminal interval, the complete intersection supplies the fixed \(1/2\) and \(3/2\) readouts.
Those readouts recover \(H\) and its first dissipative height without asking a scalar identity to manufacture their bounds.

Any finite number of temporal rows reaches only a finite spatial height.
No one finite row can contain the entire climb demanded by an alleged finite-time escape.
The object forced by this discovery is simultaneous membership in every finite temporal and spatial Sobolev row.
